% !TEX program = xelatex
% Компіляція: xelatex presentation.tex && xelatex presentation.tex
\documentclass[aspectratio=169,12pt]{beamer}

\usepackage{fontspec}
\usefonttheme{serif}
\setmainfont{Times New Roman}
\newfontfamily\cyrillicfont{Times New Roman}[Script=Cyrillic]
\newfontfamily\cyrillicfontsf{Times New Roman}[Script=Cyrillic]
\newfontfamily\cyrillicfonttt{Courier New}[Script=Cyrillic]

\usepackage{polyglossia}
\setdefaultlanguage{ukrainian}

\usepackage{graphicx}
\usepackage{array}
\usepackage{booktabs}
\usepackage{amsmath}
\usepackage{hyperref}
\usepackage{xurl}

\graphicspath{{images/}}

\hypersetup{colorlinks=true, linkcolor=black, urlcolor=blue}

\setbeamertemplate{navigation symbols}{}
\setbeamertemplate{footline}{%
  \hfill\insertframenumber/\inserttotalframenumber\kern1em\vspace{2pt}%
}

\setbeamertemplate{itemize item}{\textbullet}
\setbeamertemplate{itemize subitem}{\textendash}

\IfFileExists{images/fig-05-ui.png}{%
  \newcommand{\UiScreenshot}{fig-05-ui.png}%
}{%
  \newcommand{\UiScreenshot}{fig-03-architecture.png}%
}

\newcommand{\maybeincludegraphics}[2][0.88\textwidth]{%
  \IfFileExists{images/#2}{%
    \includegraphics[width=#1]{#2}%
  }{%
    \fbox{\parbox[c][4.5cm][c]{#1}{%
      \centering\small
      \textbf{Місце для рисунка}\\[0.5em]
      \texttt{thesis/images/#2}%
    }}%
  }%
}

\title{Монокулярна візуальна одометрія\\в реальному часі для БПЛА}
\subtitle{Курсова робота}
\author{Уточкіна Яна Максимівна}
\institute{%
  КНУ імені Тараса Шевченка\\
  Факультет комп'ютерних наук та кібернетики\\
  Кафедра математичної інформатики\\[0.4em]
  Науковий керівник: док. ф.-м. н. Терещенко В.~М.%
}
\date{2026}

\begin{document}

% --- Слайд 1 ---
\begin{frame}
  \titlepage
\end{frame}

% --- Слайд 2 ---
\begin{frame}{Актуальність}
  \begin{itemize}
    \item Автономна навігація БПЛА потребує оцінки руху без GNSS.
    \item Бортова камера - доступне джерело інформації про зміщення та обертання.
    \item \textbf{Візуальна одометрія (VO)} інтегрує відносний рух між кадрами в траєкторію.
    \item Монокулярна VO - одна камера: простіша апаратура, але неоднозначність масштабу та дрейф.
    \item Повні SLAM-системи (ORB-SLAM, PTAM) складні для навчання - потрібен доступний веб-прототип.
  \end{itemize}
\end{frame}

% --- Слайд 3 ---
\begin{frame}{Мета та завдання}
  \textbf{Мета:} розробити веб-систему монокулярної VO з інтерактивною 3D-візуалізацією.

  \medskip
  \textbf{Завдання:}
  \begin{enumerate}
    \item проаналізувати теоретичні основи VO та калібрування;
    \item спроєктувати клієнт-серверну архітектуру;
    \item реалізувати модуль оцінювання руху (оптичний потік, епіполярна геометрія);
    \item реалізувати калібрування камери за шаховою дошкою;
    \item забезпечити візуалізацію траєкторії та профілі камери;
    \item перевірити на даних EuRoC MAV (калібрування, VO, ATE за еталоном).
  \end{enumerate}
\end{frame}

% --- Слайд 4 ---
\begin{frame}{Теоретична основа}
  \begin{columns}[T,onlytextwidth]
    \column{0.48\textwidth}
    \textbf{Модель камери}
    \begin{itemize}
      \item центральна перспектива, матриця $K$;
      \item калібрування за шаховою дошкою;
      \item похибка репроєкції - критерій якості.
    \end{itemize}

    \medskip
    \textbf{Оцінювання руху}
    \begin{itemize}
      \item Shi-Tomasi + Lucas-Kanade;
      \item суттєва матриця $E$ + RANSAC;
      \item \texttt{recoverPose} $\rightarrow R$, $\mathbf{t}$.
    \end{itemize}

    \column{0.48\textwidth}
    \centering
    \maybeincludegraphics[0.95\columnwidth]{fig-01-pinhole.png}\\[0.3em]
    {\footnotesize Модель камери з центральною перспективою}
  \end{columns}
\end{frame}

% --- Слайд 5 ---
\begin{frame}{Відстеження ознак}
  \centering
  \maybeincludegraphics[0.72\textwidth]{fig-02-optical-flow.png}

  \medskip
  {\small
  Зелені - ознаки на кадрі $t{-}1$;\;
  сині - успішно відстежені;\;
  червоні - відхилені (forward-backward перевірка).
  }
\end{frame}

% --- Слайд 6 ---
\begin{frame}{Архітектура системи}
  \centering
  \maybeincludegraphics[0.82\textwidth]{fig-03-architecture.png}

  \medskip
  \begin{itemize}
    \item \textbf{Сервер} (Python, FastAPI, OpenCV) --- візуальна одометрія та калібрування.
    \item \textbf{Клієнт} (Angular, Plotly.js) - захоплення відео, 3D-графік.
    \item \textbf{WebSocket} - обмін кадрами (JPEG/Base64) і координатами в реальному часі.
  \end{itemize}
\end{frame}

% --- Слайд 7 ---
\begin{frame}{Конвеєр візуальної одометрії}
  \centering
  \maybeincludegraphics[0.78\textwidth]{fig-04-vo-pipeline.png}

  \medskip
  {\small
  Резервний шлях - афінна оцінка (\texttt{estimateAffinePartial2D});\;
  масштаб: тріангуляція + оптичний потік + базова лінія ключових кадрів;\;
  прив'язка пози до ключового кадру кожні 45~кадрів.
  }
\end{frame}

% --- Слайд 8 ---
\begin{frame}{Технологічний стек}
  \small
  \begin{table}
    \centering
    \begin{tabular}{@{}>{\raggedright\arraybackslash}p{2.4cm}>{\raggedright\arraybackslash}p{8.8cm}@{}}
      \toprule
      Шар & Технології \\
      \midrule
      Сервер & Python 3.11, FastAPI, OpenCV, NumPy, PyYAML \\
      Клієнт & Angular 21, TypeScript, Angular Material, Plotly.js \\
      Транспорт & WebSocket, REST (\texttt{/api/configs/calibrate}) \\
      Розгортання & Docker Compose (порти 8000 / 4200) \\
      \bottomrule
    \end{tabular}
  \end{table}

  \medskip
  Режими: \textbf{«Потік»} (вебкамера) та \textbf{«З файлу»} (відео EuRoC).
  Профілі камери - YAML у \texttt{vo-uav/storage/configs/}.
\end{frame}

% --- Слайд 9 ---
\begin{frame}{Результати калібрування (EuRoC)}
  \begin{columns}[T,onlytextwidth]
    \column{0.44\textwidth}
    \small
    \begin{table}
      \centering
      \begin{tabular}{@{}lr@{}}
        \toprule
        Параметр & Значення \\
        \midrule
        $f_u$ & 397{,}88~px \\
        $f_v$ & 394{,}66~px \\
        $c_u$ & 360{,}80~px \\
        $c_v$ & 271{,}66~px \\
        RMS репроєкції & \textbf{1{,}21~px} \\
        Сітка & 6 $\times$ 7 \\
        \bottomrule
      \end{tabular}
    \end{table}

    \vspace{0.4em}
    1740 з 2376 кадрів \texttt{cam\_checkerboard}  детекція дошки.

    \column{0.52\textwidth}
    \centering
    \maybeincludegraphics[0.98\columnwidth]{fig-06-calibration.png}
  \end{columns}
\end{frame}

% --- Слайд 10 ---
\begin{frame}{Результати VO на EuRoC \texttt{V1\_01\_easy}}
  \small
  \begin{columns}[T,onlytextwidth]
    \column{0.48\textwidth}
    \begin{table}
      \centering
      \begin{tabular}{@{}lr@{}}
        \toprule
        Метрика & Значення \\
        \midrule
        Успішних кадрів & 99,9\,\% \\
        ATE (СКП) & \textbf{1,74~м} \\
        $s$ (Sim(3)) & 0,061 \\
        Довжина еталону & 58,6~м \\
        Шлях VO (масштаб.) & \textbf{24,1~м} \\
        \bottomrule
      \end{tabular}
    \end{table}

    \vspace{0.5em}
    \textbf{24~м}
    ($\approx$41\,\% довжини еталону), ATE $\approx$1,7~м.

    \column{0.52\textwidth}
    \centering
    \maybeincludegraphics[0.98\columnwidth]{vo-eval-V1_01_easy.png}
  \end{columns}
\end{frame}

% --- Слайд 11 ---
\begin{frame}{Наочна демонстрація}


\end{frame}

% --- Слайд 12 ---
\begin{frame}{Обмеження та перспективи}
  \begin{columns}[T,onlytextwidth]
    \column{0.48\textwidth}
    \textbf{Обмеження}
    \begin{itemize}
      \item дрейф без замикання петлі;
      \item умовний масштаб під час польоту (метри - лише з еталону);
      \item залежність від якості калібрування;
      \item затримка WebSocket у режимі «Потік».
    \end{itemize}

    \column{0.48\textwidth}
    \textbf{Перспективи}
    \begin{itemize}
      \item відносна помилка траєкторії (RPE);
      \item ORB-дескриптори, локальне пакетне уточнення;
      \item поєднання з IMU (VIO);
      \item замикання петлі.
    \end{itemize}
  \end{columns}
\end{frame}

% --- Слайд 13 ---
\begin{frame}{Висновки}
  \begin{itemize}
    \item Реалізовано \textbf{веб-систему монокулярної VO} з клієнт-серверною архітектурою.
    \item Модуль VO: Shi-Tomasi, Lucas-Kanade, $E$ + RANSAC, афінний резерв, прив'язка до ключових кадрів.
    \item Калібрування за шаховою дошкою з автопідбором сітки та профілями YAML.
    \item На EuRoC \texttt{V1\_01\_easy}: RMS \textbf{1,21~px}; $s \approx 0{,}061$; ATE $\approx 1{,}7$~м; шлях VO $\approx$24~м (еталон 58,6~м).
    \item Система придатна для \textbf{навчальної демонстрації} та прототипування локалізації БПЛА.
  \end{itemize}
\end{frame}

% --- Слайд 14 ---
\begin{frame}
  \centering
  \Large
  \textbf{Дякую за увагу!}\\[1.2em]
  Готова відповісти на запитання.
\end{frame}

\end{document}
